\begin{textsample}{NEG Dim 2 – generic\_gpt – Score -3.00 – generic\_gpt/t1750\_gpt.txt}  \label{ex:f2_neg_017}
I ’ve been feeling this weird \textbf{mix} of invisible and exposed lately, and I don't really know how to explain it to anyone in my life without sounding dramatic, so I ’m putting it here instead.
On paper I ’m “ not alone. ” I have coworkers, a couple of old friends I text now and then, family chats that light up on holidays.
But it all feels like background \textbf{noise}.
I can go days without anyone asking me a real question about how I ’m doing, and if they do, I give the safe answer because I ’m not sure they actually want to know.
I think what ’s getting to me is the sense that I ’m just not really *known* by anyone. \textbf{People} know facts about me—my job, my hobbies, where I live—but not the stuff that actually keeps me awake at night or the small things that make me feel like a person.
I ’ll have a \textbf{conversation}, laugh at the right places, say the right things, and then go home feeling like I played a \textbf{role} instead of being myself.
I used to think loneliness was about being physically alone, but lately it feels more like this constant mismatch between the amount of connection I need and the amount I actually have.
I scroll through social media, see everyone ’s group photos and inside jokes, and it ’s like watching life from the outside of a \textbf{window}.
I ’m there, but not really “ in ” anything.
I don't even know what I ’m expecting from posting this.
I guess I just wanted to put this somewhere, to say out loud ( or as close as I can get ) that I ’m tired of feeling like a background character in my own life.
It ’s starting to feel less like a phase and more like the default setting.

% matched lemmas: conversation, mix, noise, people, role, window
\end{textsample}
